\section{Discussion}\label{sec:discussion}

This section discusses the implications of the Phase I results for theoretical physics, compares the SS-PEG framework with existing approaches, and outlines the path toward Papers II and III.

\subsection{Comparison with Existing Literature}\label{sec:comparison_lit}

\subsubsection{Connes' Spectral Action Principle}

Connes' spectral action principle \cite{Connes1994,ConnesLott1990,ChamseddineConnes1997} comes closest to our approach: it derives the Standard Model Lagrangian from the spectral triple of a non-commutative geometry on $M_4 \times F$. However, Connes \textit{chooses} the finite space $F$ to reproduce the SM gauge group. We claim the SM is the \textbf{generic} outcome of spectral optimization without such a choice.

The key distinction is:

\begin{itemize}
\item \textbf{Connes}: The Cartan--Weyl graph is the \textit{output}---chosen to match observation.
\item \textbf{SS-PEG}: The Cartan--Weyl graph is the \textit{input}---it emerges from the variational principle.
\end{itemize}

This inversion changes the logical direction: Connes starts from the SM and derives its Lagrangian; we start from a geometric principle and derive the SM.

\subsubsection{Emergent Symmetry vs.\ Imposed Symmetry}

The standard paradigm in gauge theory is:

\begin{enumerate}
\item Postulate a symmetry group $G$ (e.g., SU(3) $\times$ SU(2) $\times$ U(1)).
\item Derive field equations from gauge invariance.
\item Fit parameters (coupling constants, masses).
\end{enumerate}

The SS-PEG results invert this entirely:

\begin{enumerate}
\item Specify relational density and geometric constraints.
\item The gauge group \textit{emerges}.
\item All invariants are \textit{predictions}.
\end{enumerate}

This is not merely a change of formalism. It changes what counts as an ``explanation.'' In the old paradigm, the question ``why SU(3)?'' has no answer---it is an axiom. In the new paradigm, the answer is: ``because the relational graph has $d^2 - 1 = 8$ independent generators in a 3-dimensional state space, and the Killing metric drives them to close.''

\subsection{Theoretical Implications}\label{sec:implications}

\subsubsection{Gravity and Gauge Symmetry Share a Common Origin}

The so(3,1) result shows that the same variational principle that generates internal gauge symmetries also generates spacetime symmetry---with the Killing signature encoding the causal structure. Section~\ref{sec:spacetime} establishes that gravity and gauge forces are \textit{two faces of the same Killing form}: the compact eigenvalues produce gauge forces, the non-compact eigenvalues produce spacetime geometry. Both emerge simultaneously at the density threshold $\rho_c \approx 1$.

The P$\to$R$\to$C$\to$I mapping (Section~\ref{sec:gravity_prci}) demonstrates that the relational chain unifies gauge theory and general relativity at the structural level. The Riemann curvature tensor arises as the cycle discrepancy of the Levi-Civita connection in exactly the same way that the gauge field strength arises as the cycle discrepancy of the gauge connection. The difference is only in the base group: SO(3,1) for gravity, SU(N) for gauge forces.

Following Jacobson~\cite{Jacobson1995,Jacobson1998}, the Einstein equations emerge from entanglement thermodynamics $\delta S = \delta\langle K \rangle$. The relational density $\rho$ provides the microscopic substrate: high $\rho$ generates entanglement entropy, gradients of $\rho$ produce the modular Hamiltonian, which yields effective acceleration and local curvature via the Unruh effect. Gravity is not fundamental but the thermodynamic equilibrium condition of the entanglement structure in the evolution graph.

\subsubsection{Dark Matter as Connection Without Particles}

Section~\ref{sec:dark_matter} proposes that ``dark matter'' is not a particle but the footprint of the Levi-Civita connection in regions of intermediate relational density. This yields distinctive predictions---filamentary/membrane topology, cored density profiles, no self-annihilation, no radiation---that are testable via weak lensing surveys. Recent observations of ultrafaint dwarf galaxies~\cite{SanchezAlmeitaTrujilloPlastino2024} (cored profiles, universal radial profile, feedback excluded) provide early evidence consistent with the connection-without-particles picture and disfavor collisionless CDM cusps at $\geq$97\% confidence.

\subsubsection{Symmetry Is Emergent, Not Fundamental}

The deepest lesson of the experimental program is that \textbf{symmetry is not a fundamental ingredient of physics but an emergent consequence of relational geometry}. This inverts a century of theoretical tradition:

\begin{itemize}
\item \textbf{Commutation relations are not fundamental.} They are derived quantities, emergent from the Killing geometry. Test 3.2 proves the converse: imposing commutation structure \textit{without} Killing geometry produces nothing.

\item \textbf{The Jacobi identity is automatic.} Within matrix representations (the natural setting for quantum mechanics), $J(X,Y,Z) = 0$ holds for \textit{any} set of matrices, not just Lie algebras. It is an algebraic triviality, not a physical constraint. The real constraint is \textit{closure}---whether commutators remain within the generator span.

\item \textbf{Representation theory is encoded in the Killing metric.} The dual Coxeter number $h^\vee$, the Dynkin index $I_R$, and the Casimir eigenvalues are all readable directly from the emerged Killing eigenvalues via $h^\vee = I_R \cdot |K|/\kappa^2$. This formula works for 15/15 algebras with zero exceptions.
\end{itemize}

\subsubsection{Spacetime Is Algebraic, Not an Arena}

The so(3,1) result shows that the same variational principle that generates internal gauge symmetries also generates spacetime symmetry---with the Killing signature encoding the causal structure. This provides a path toward understanding why spacetime has Lorentzian (not Euclidean) signature:

\begin{quote}
The Minkowski metric $\eta_{\mu\nu} = \mathrm{diag}(-1,+1,+1,+1)$ emerges because the Lorentz algebra so(3,1) has Killing signature $(3+,3-)$, which encodes exactly 3 compact (spatial rotation) and 3 non-compact (boost) directions.
\end{quote}

The ``3+1'' dimensionality of spacetime is thus related to the representation theory of the Lorentz group: 6 generators splitting into 3 rotations and 3 boosts in the 4-dimensional representation.

\subsubsection{The Number of Free Parameters May Be Zero}

The framework proposes that the three coupling constants of the Standard Model are not free parameters but emergent properties of the Cartan--Weyl graph's spectral structure. Paper III derives $\alpha_{\mathrm{EM}}$, $\alpha_s$, and $\alpha_W$ from the Ramanujan ratios of the SU(2) and SU(3) subgraphs with sub-percent precision. If confirmed, this would mean the number of free parameters in the Standard Model is \textbf{zero}---all observable quantities are predictions of a single geometric principle.

\subsection{Limitations and Open Questions}\label{sec:limitations}

\subsubsection{Continuum Limit ($d \to \infty$)}\label{sec:continuum_limit}

All current experiments work with finite-dimensional matrix representations. The connection to continuum field theory requires taking $d \to \infty$ in a controlled way, ideally showing that the gauge field equations of motion emerge in this limit. This is a critical open question for the program.

\subsubsection{Connecting Phase I with Phase II--III}\label{sec:connecting_phases}

Phase I establishes the emergence of Lie algebras from the Killing metric. Phase II (Paper II) develops the graph-topological interpretation, classifying edges as cyclic, free, or decoupled according to the Killing spectrum, and establishing the Ramanujan property as a criterion of spectral optimality. Phase III (Paper III) derives the Standard Model mass spectrum, coupling constants, and flavor mixing from the same graph topology with sub-percent precision.

The bridge between Phase I and Phase II is the Cartan--Weyl graph: the discrete topological object that encodes the algebraic structure in a form amenable to spectral analysis. The bridge between Phase II and Phase III is the master formula for coupling constants: $\alpha_X = e^{S_X}/(4\pi)$ with direction vectors $\mathbf{n}_X$ read from the cycle structure of the CW graph.

\subsubsection{Deriving the Five Building Blocks from a Single Principle}\label{sec:building_blocks}

Paper III identifies five algebraic invariants from the Cartan--Weyl root graph of the Standard Model gauge algebra $\su(3) \oplus \su(2) \oplus \fu(1)$:

\begin{equation}
\begin{array}{ll}
\mathcal{R}_2 = 1/2 & \text{Spectral ratio of SU(2) subgraph} \\
\mathcal{R}_3 = \frac{\sqrt{57}-3}{2\sqrt{17}} \approx 0.5523 & \text{Spectral ratio of SU(3) subgraph} \\
R_{\mathrm{EE}} = 1/\sqrt{2} \approx 0.7071 & \text{Euler--Euler spectral ratio (cross-talk)} \\
h^\vee_3 = 3 & \text{Dual Coxeter number of SU(3)} \\
h^\vee_2 = 2 & \text{Dual Coxeter number of SU(2)}
\end{array}
\end{equation}

These five numbers are \textit{not free parameters}---they are fixed by the topology of the root system. Deriving them from a single geometric principle (rather than reading them off the CW graph) remains an open question.

\subsection{Connection with Papers II and III}\label{sec:connection_papers}

\subsubsection{Paper II: Spectral Quality of the Emergent Killing Form}

Paper II develops the spectral analysis of the Cartan--Weyl graphs associated with each emergent algebra. Key results include:

\begin{itemize}
\item The SM factors SU(2) and SU(3) are \textbf{deeply Ramanujan} (ratios 0.50 and 0.55), passing all six alternative definitions unanimously.
\item A sharp spectral phase transition occurs at Coxeter number $h^* \approx 9.5$: algebras with $h < h^*$ are Ramanujan (ordered phase); those with $h > h^*$ are non-Ramanujan (disordered phase).
\item The SM lives deep inside the \textbf{ordered phase}.
\item Ramanujan $\times$ Ramanujan $\to$ Ramanujan: the SM product $\su(3) \oplus \su(2)$ is unanimously Ramanujan (ratio 0.657).
\end{itemize}

\subsubsection{Paper III: Standard Model from Graph Topology}

Paper III derives the Standard Model mass spectrum, coupling constants, and flavor mixing from the Cartan--Weyl graph topology. Key results include:

\begin{itemize}
\item The three coupling constants emerge from the Ramanujan ratios of the SU(2) and SU(3) subgraphs with sub-percent precision.
\item Mass ratios are predicted from the five building blocks with mean absolute error 0.32\%.
\item CKM and PMNS mixing angles are derived from the graph structure with $> 95\%$ of targets within 1\% of experimental values.
\end{itemize}

The three papers form a logical chain: Phase I establishes \textit{which} algebras emerge; Phase II establishes \textit{why} the SM is spectrally optimal among them; Phase III derives \textit{what} observable quantities the SM graph predicts.
